\section{What happened, in order}
\label{sec:story}

This section is the history of the project. It explains why the system looks
the way it does: most design decisions were forced by something that went
wrong on a real server.

\subsection{Before meeting 1 (until 16 October)}

\paragraph{Stage 1: the harness.} An executor (\file{src/pgfuzz/reference/pg_executor.py})
runs a test case against a real \pg server, one savepoint per statement,
rolled back at the end; it reads the server log for \texttt{TRAP} (assertion),
\texttt{PANIC} and ``terminated by signal'' lines; it recovers after crashes.
An oracle (\file{src/pgfuzz/eval/oracle.py}) classifies each test case
(crash, assertion, internal error \texttt{XX000}, timeout, \dots). A first
wrong-result oracle re-ran each deterministic query under 17 planner settings
(the \emph{plan-variant} oracle). Every trial writes a resumable run log
(\file{events.jsonl}, \file{timeline.jsonl}, \file{summary.json}) and a
report builder produces tables, plots and statistics.

\paragraph{Stage 1: the random/mutation arm} (\file{src/generator/}). Seeds
are built with a small SQL AST; they are then mutated as \emph{text} (flip a
character, delete a span, swap an operator, splice a keyword). The mutators
know nothing about types or scopes, so many mutants are invalid. It ignores
feedback on purpose: it is the unguided baseline.

\paragraph{Stage 2: the typed grammar.} We designed our own grammar format
(\texttt{.pgg}) and a derivation engine with scopes and type variables
(Section~\ref{sec:grammar}). Four tuning rounds of 300 test cases each
on the same server and seed raised server-side validity:

\begin{center}
\begin{tabular}{lrrrrrr}
\toprule
Round & Statements & Accepted & Semantic err. & Runtime err. & Timeouts & Time \\
\midrule
1 & 812 & 677 (83.4\%) & 76 & 57 & 0 & 16.7 s \\
2 & 828 & 688 (83.1\%) & 39 & 89 & 7 & 82.9 s \\
3 & 929 & 821 (88.4\%) & 28 & 74 & 6 & 56.1 s \\
4 & 966 & 966 (100\%) & 0 & 0 & 0 & 22.1 s \\
\bottomrule
\end{tabular}
\end{center}

Round 2 looks like no progress (83.1\%), but the mix changed: semantic errors
halved, more statements reached execution, and one rule (\texttt{generate\_series}
up to a table column) wasted the clock with million-row series. The edits
were always driven by the error list, never at random. This is the number we
showed at meeting 1. It is a validity number, not a bug.

\subsection{Meeting 1 (16 October)}

We presented the three claims, the architecture, the RQs, and a plan with
seven priorities: (1) long grammar campaigns, (2) metamorphic oracles on
derivation trees, (3) cross-version testing, (4) triage automation, (5) the
commit-directed LLM loop with diff coverage, (6) a re-injected bug benchmark
with statistics, (7) stretch goals (plan-shape feedback, a CERT-style
performance oracle, a row-level security oracle).

\subsection{The bug hunt (16--17 October)}

\paragraph{Setup.} In a cloud container we built \pg master at
\texttt{26bd08fb} with assertions and \pg 18.6 (\texttt{REL\_18\_STABLE}) as a
reference. We ran the grammar campaign (\file{configs/cloud_hunt.yaml}, two
trials of six hours) and two cross-version join fuzzers
(\file{src/xversion/joinfuzz*.py}) that run the same query on master and on
18.6 with random planner settings.

\paragraph{Candidates.}
\begin{enumerate}
  \item A plan-variant wrong result under \texttt{enable\_hashagg=off}:
    \emph{false positive}. \texttt{ARRAY(SELECT \dots{} EXCEPT ALL \dots)} has no
    \texttt{ORDER BY}, so the array's element order legitimately depends on the
    plan. Fixed in the determinism filter \texttt{eligible()}.
  \item An assertion failure about 600 s into trial 0 (test case 13306):
    \texttt{TRAP: failed Assert("list\_length(rcexpr->inputcollids) == nopers")},
    \file{execExpr.c:2078}. \textbf{A real bug (BUG-01)}, see
    Section~\ref{sec:bug01}.
  \item More \texttt{ARRAY(SELECT)} false positives from a trial that started
    before the fix.
  \item A crash-recovery failure on one server after the BUG-01 crash chain:
    \texttt{FATAL: could not create file "base/5/165792": File exists}. We traced
    it in the WAL to an unlogged sequence created in an open transaction, read
    the recovery code (\texttt{ResetUnloggedRelations}), but could not
    reproduce it from a clean data directory. Unreproduced bugs are not
    reported. Evidence is kept in \file{docs/bughunt/evidence/}.
\end{enumerate}

\paragraph{Lessons that became code.} The failure writer once saved the wrong
test case for a signature; the determinism filter needed the
\texttt{ARRAY(SELECT)} rule; a cross-version mismatch came from
\texttt{DEFAULT now()} in our own seed data (the clock differs between
servers) -- fixed with deterministic data and by excluding clock functions.

\subsection{The build session (23 October)}

This is where most of the meeting-2 system was built, in this order.

\begin{enumerate}
  \item \textbf{Typed node spans.} The engine wraps each node's tokens in
    invisible markers, so after rendering it knows exactly which tokens each
    tree node produced. This is what lets an oracle edit one node.
  \item \textbf{Metamorphic oracles on derivation trees} (TLP, NoREC,
    expression rewrites, CTE toggle, table fence). In the very first sample
    of 150 test cases, the scalar-subquery rewrite \texttt{e} $\rightarrow$
    \texttt{(SELECT e)} found \textbf{BUG-02} (Section~\ref{sec:bug02}).
    A 600-test-case live check then gave 907 checks, 0 false positives.
  \item \textbf{Plan-shape guidance and metric.} A test case whose plan has a
    new shape (digest of node types in \texttt{EXPLAIN}) enters the corpus.
  \item \textbf{Cross-version executor} (master vs 18.6) with a sequence
    prelude, deterministic data, clock functions excluded: 1{,}155 comparisons,
    0 mismatches after the fixes.
  \item \textbf{Triage pipeline}: reproduce, reduce (statements, clauses,
    tokens; pairs for metamorphic checks), report, bisect script, precision
    ledger. On BUG-02's candidate: 4 statements / 107 tokens $\rightarrow$
    1 statement / 36 tokens in 307 runs.
  \item \textbf{LLM arm}: commit miner, cached prompts, client with manual
    mode, validation gate, diff coverage, repair loop, ledger.
  \item \textbf{RQ3 benchmark}: 23 fixes reverted. The first PoC verification
    showed only 12/23 observable. Two of our own mistakes caused most of the
    gap: the PoC tables were called \texttt{i4}/\texttt{i8}, names our schema
    already uses for \emph{indexes}, so the PoCs failed at \texttt{CREATE TABLE};
    and the digest masked float aggregates so completely that \texttt{0} vs
    \texttt{NaN} looked equal. After renaming (prefix \texttt{rb\_}), keeping
    NaN/Infinity in masked floats, adding \texttt{ANALYZE}, and fixing two PoCs:
    \textbf{20 of 23 observable}, the same symptom with all 23 injected
    together (no masking).
  \item \textbf{Smoke run 1} (4 min per arm on the injected server): the
    grammar arm \textbf{rediscovered BUG-01}; five metamorphic candidates turned
    out to be \textbf{our} false positives: TLP-HAVING removed \texttt{HAVING}
    from a level without \texttt{GROUP BY} (turning a one-group aggregate into a
    plain query), and \texttt{DISTINCT} over \texttt{numeric} may keep
    \texttt{21.0000} or \texttt{21.00} depending on the plan. Both fixed with
    tests. Also: all four workers shared one reference server and raced on
    schema setup; fixed with one database per worker.
  \item \textbf{Smoke run 2} (15 min per arm): Grammar+Oracles and +LLM found
    re-injected bug \texttt{rb-dfcf6fa7} at about 160 s; attribution by replay
    on the solo servers picked the right one of two bugs that share an error
    message. Two random-arm ``finds'' were \textbf{our} fault: a PoC replay we
    ran during the campaign crashed the shared server and the harness blamed
    the random arm's statement. Fixed: crash evidence now must carry the test
    case's own backend PID. A cross-version difference (\texttt{CORR} of
    constant input: NULL on master, 1 on 18.6) is an intended master change;
    those functions are excluded.
  \item \textbf{Paper draft, execution plan, progress report.}
\end{enumerate}

\subsection{The plan-coverage round (24 October)}

We re-read the meeting-1 plan and slides line by line and implemented every
item still missing: TLP for aggregates and the remaining rewrites of plan
Section~3.4; grammar-tree reduction (``one tree, four jobs''); localization of
wrong results (plan oracle $\times$ rewrite oracle); feedback-weighted grammar
alternatives (bandit); a data profile for boundary and selective literals;
the RLS leak oracle; the CERT oracle; LLM-proposed equivalence rewrites with a
value probe; the eager-aggregation plan variant; upstream bug-report drafts;
95\% bootstrap confidence intervals. Each was checked on master
(Section~\ref{sec:components}). The traceability matrix is
\file{docs/PLAN_COVERAGE.md}.

\begin{warnbox}[Things that went wrong and what they taught us]
\begin{itemize}
  \item A richer grammar exposes more corners where two ``equivalent'' queries
    are not equivalent; oracle precision, not generator power, limits a
    wrong-result fuzzer. Every false-positive class became a fix plus a test.
  \item Shared servers are dangerous: our own manual replay corrupted a
    campaign's attribution; the fix is PID-checked evidence.
  \item Our own PoCs were wrong twice (name clash, stats-dependent plans);
    ground truth needs verification too (\texttt{bench verify}).
  \item Sampling is not proof: a false LLM rewrite passed 109 sampled checks
    and was caught only by evaluating it on NULL.
  \item \texttt{pkill -f pattern} from a shell whose command line contains the
    pattern kills that shell; wait loops that \texttt{pgrep} their own pattern
    never end. Small, but they cost us runs.
\end{itemize}
\end{warnbox}
